\section{Практические ограничения и направления развития}
\label{sec:limitations}

\subsection{Ограничения текущей реализации}

Несмотря на положительные результаты по линии «микросервис + потоковое
обновление», текущая реализация имеет ряд ограничений.

Во-первых, \texttt{Dump} остаётся дороже, чем у \texttt{ranger}. В прикладном
сценарии это допустимо, если полный снимок вызывается сравнительно
редко, но при интенсивном использовании полного дампа может
понадобиться дополнительный кэшированный слой полного снимка или
потоковая выдача результата. Основная стоимость связана с построением
непересекающейся таблицы для внешних потребителей.

Во-вторых, полный \texttt{Dump} на миллионах записей создаёт большой
поток аллокаций. При жёстком \texttt{GOMEMLIMIT} операция завершается,
но замедляется из-за работы GC, поэтому production-конфигурация должна
учитывать параллельную нагрузку, gRPC buffers и возможную ресинхронизацию.

В-третьих, потоковая модель пока в основном рассматривается как
последовательность update-событий из одного источника. Для более
сложных промышленных сценариев могут понадобиться расширенные гарантии
повторной доставки, дедупликации и диагностики пропущенных событий
\cite{kreps2011kafka,kleppmann2017ddia}.

\subsection{Возможные доработки реализации}

Ограничения, выявленные в экспериментальной части, указывают на три
группы доработок текущей реализации.

\textit{Lookup.} Операция LPM поверх ART использует наборы активных длин
префикса и поэтому проверяет только релевантные кандидатные маски.
Дальнейшая доработка может состоять в хранении признака терминальности
непосредственно в узлах дерева, сохранении последнего найденного
валидного совпадения во время прохода по ключу и разделении быстрых путей
обработки для IPv4 и IPv6. Работы
Tree~Bitmap~\cite{eatherton2004treebitmap} и
Poptrie~\cite{asai2015poptrie} показывают близкие приёмы оптимизации
поиска наиболее специфичного префикса.

\textit{Dump.} Для редких, но дорогих полных выгрузок логично ввести
кэшированный упорядоченный снимок, ленивую материализацию или
потоковую выдачу содержимого вместо одного большого ответа. Если
сервис начнёт использоваться как поставщик периодических выгрузок,
полезны и дифференциальные снимки для конкретных подписчиков.

\textit{Телеметрия.} Расширение телеметрии и интеграция с
BMP-наблюдением~\cite{rfc7854} и публичными
коллекторами~\cite{ripe-ris,routeviews} позволят измерять задержку
актуализации, число принятых и пропущенных update-событий, время
ресинхронизации, размеры очередей и ошибки соединения с GoBGP. Сравнение
локальной таблицы с независимым внешним источником также может служить
сигналом потери или неправильной интерпретации событий. Здесь же уместен
механизм, аналогичный route flap damping~\cite{rfc2439}, для подавления
гиперактивных префиксов.

Шаблон «отдельный сервис с единым API + снимок + поток + шардированное
in-memory ядро» переносим за пределы рассмотренной пары продуктов. Он
применим в системах сетевой безопасности, аналитике трафика,
операторской отчётности и сетевом инвентаре -- везде, где нужно держать
актуальный маршрутный контекст для принятия решений и не дублировать
логику между потребителями.

%==============================================================
